% ==================================================================================================================================
% Introduction

\minitoc  % Affiche la table des matières pour ce chapitre

Dans ce chapitre, nous allons définir formellement les matrices ainsi que les espaces vectoriels de matrices. 
Ces objets vont nous permettre de représenter très facilement les applications linéaires introduites dans le 
chapitre 1. Les propriétés que nous allons voir sur les matrices pourront, pour la plupart se transposer aux applications 
représentées. 

\vspace{0.3cm}

Dans tout ce chapitre, on nommera $E$ un $\K$-espace vectoriel où $\K$ est un corps commutatif. 

% ==================================================================================================================================
% Matrices 

\section{Matrices}

\begin{definition}[Matrice]
    On appelle matrice $n \in \N$ lignes et $p \in \N$ colonnes toute application 
        \[ M : \{1, \dots , n\} \times \{1, \dots , p\} \longrightarrow \K \] 
    On dit alors que $M$ est de taille $n + p$. 
\end{definition}

\begin{example}
    Les matrices sont représentées de façon tabulaire. 
    Par exemple une matrice $M$ à 2 lignes et 3 colonnes à coefficients dans $\R$ peut être : 
        \[ M = 
            \begin{pmatrix}
                1 & 26 & \frac{3}{2} \\ 
                \pi & 6 & \sqrt{2} 
            \end{pmatrix} \] 
    On notera $M(i,j)$ ou $M_{i,j}$ le coefficient de $M$ de la i-ème ligne et j-ème colonne. 
    Ici, on a donc $M(1,1) = 1$ et $M(2, 3) = \sqrt{2}$. 
\end{example}

\begin{proposition}[Matrices Remarquables]
    Le nombre de lignes et de colonnes d'une matrice pouvant varier dans $\N \times \N$, il existe des 
    matrices dites remarquables : 
    \begin{itemize}
        \item \textbf{Matrice ligne} : une matrice de taille $ 1 \times p$ de la forme : $(x_1, \dots, x_p)$ 
        \item \textbf{Matrice colonne} : de la forme 
        $ \begin{pmatrix}
            x_1 \\ 
            \vdots \\ 
            x_p 
        \end{pmatrix}$ 
    \end{itemize} 
\end{proposition}

\begin{proposition}[Matrice Égales]
    On dit que deux matrices M et N $ \in \mathcal{M}_{n,m}(\K)$ sont égales ssi elles sont de même taille et 
    tout leurs coefficients sont égaux. 
        \[ \text{i.e} \quad \forall i,j \in \llbracket 1, n \rrbracket \times \llbracket 1, m \rrbracket, M(i,j) = N(i,j) \]  
\end{proposition}

\begin{definition}[Transposition]
    Soit $ \mathcal{M}_{n,m}(\K)$ l'ensemble des matrices à coefficients dans un corps $\K$. 
    On défini l'application transposition comme : 
        \[ T : 
            \begin{cases}
                \mathcal{M}_{n,m}(\K) \longrightarrow \mathcal{M}_{m,n}(\K) \\ 
                M = (m_{i,j}) \longrightarrow ^t M = (m_{j,i})
            \end{cases} \]
    Cette application prend une matrice $ n \times m$ et la retourne en une matrice $ m \times n $ en "transformant"
    ses lignes en colonnes. 
    Elle est \textbf{involutive}, autrement dit $ \forall M \in \mathcal{M}_{n,m}(\K) \; ^t (^t M) = M$ 
    et \textbf{linéaire}. 
\end{definition}

\begin{definition}[Matrice symétrique]
    On dit qu'une matrice $M \in \mathcal{M}_{n}(\K)$ est symétrique ssi $^tM = M$. 
\end{definition}


% ==================================================================================================================================
% Structure 

\section{Structure}

L'ensemble $ \mathcal{M}_{n,m}(\K)$ des matrices de taille $ n \times m $ à coefficients dans $\K$ est 
un espace vectoriel de dimension $ n \times n$. 